\documentclass[12pt]{beamer}
\usetheme{Madrid}
\usecolortheme{seahorse}

\usepackage[utf8]{inputenc}
\usepackage{graphicx}
\usepackage{booktabs}
\usepackage{longtable}
\usepackage{hyperref}
\usepackage{geometry}
\usepackage{amsmath}
\usepackage{caption}

\title{Tema 6 – Critérios de Elegibilidade para Revisão Sistemática}
\author{Departamento de Física, Universidade Eduardo Mondlane}
\date{\today}

\begin{document}

%-----------------------------------
\begin{frame}
\titlepage
\end{frame}

%-----------------------------------
\begin{frame}{Conteúdo}
\tableofcontents
\end{frame}

%-----------------------------------
\section{Introdução}
\begin{frame}{Introdução}
A definição de critérios de elegibilidade é uma etapa crucial na revisão sistemática, garantindo que apenas estudos relevantes sejam incluídos. Esta aula aborda os conceitos, exemplos, refinamento e ligação dos critérios à estrutura PICOTS.
\end{frame}

\begin{frame}{Objetivos de Aprendizagem}
\begin{itemize}
    \item Compreender o conceito de critérios de elegibilidade.
    \item Diferenciar critérios amplos e restritos.
    \item Identificar tipos de viés relacionados à seleção de estudos.
    \item Aplicar os critérios aos elementos do PICOTS.
\end{itemize}
\end{frame}

%-----------------------------------
\section{Visão Geral do Processo de Revisão Sistemática}
\begin{frame}{Visão Geral do Processo de Revisão Sistemática}
\begin{figure}
\centering
\includegraphics[width=0.7\textwidth]{processo_rsm.png} % Placeholder
\caption{Visão geral do processo de revisão sistemática}
\end{figure}
\end{frame}

%-----------------------------------
\section{Critérios de Elegibilidade}
\begin{frame}{Definição}
Critérios de elegibilidade determinam quais estudos serão incluídos ou excluídos da revisão, garantindo consistência e validade dos resultados.
\end{frame}

\begin{frame}{Exemplos de Critérios}
\begin{itemize}
    \item Tipo de estudo (ensaios clínicos, observacionais, revisões, etc.)
    \item População ou participantes
    \item Intervenção ou exposição
    \item Comparadores
    \item Desfecho
    \item Tempo de seguimento
    \item Idioma e ano de publicação
\end{itemize}
\end{frame}

\begin{frame}{Critérios Amplos vs. Restritos}
\begin{itemize}
    \item \textbf{Amplos:} Inclusão de mais estudos, maior heterogeneidade, menor risco de omitir dados relevantes.
    \item \textbf{Restritos:} Maior homogeneidade, resultados mais consistentes, risco de perder evidências importantes.
\end{itemize}

\begin{figure}
\centering
\includegraphics[width=0.6\textwidth]{criterios_amplos_restritos.png} % Placeholder
\caption{Exemplo de critérios amplos e restritos}
\end{figure}
\end{frame}

%-----------------------------------
\section{Viés e Seleção de Estudos}
\begin{frame}{Viés}
\begin{itemize}
    \item Erros sistemáticos na seleção de estudos podem comprometer a validade da revisão.
    \item Exemplos: viés de publicação, viés de idioma, viés de tempo.
\end{itemize}
\end{frame}

\begin{frame}{Ligação ao PICOTS}
Critérios de elegibilidade devem ser claramente vinculados aos elementos do PICOTS:
\begin{itemize}
    \item P – População
    \item I – Intervenção
    \item C – Comparação
    \item O – Desfecho
    \item T – Tempo de seguimento
    \item S – Tipo de estudo
\end{itemize}

\begin{figure}
\centering
\includegraphics[width=0.7\textwidth]{picots_criterios.png} % Placeholder
\caption{Ligação dos critérios ao PICOTS}
\end{figure}
\end{frame}

%-----------------------------------
\section{Tipos de Estudos e Considerações Adicionais}
\begin{frame}{Tipos de Estudos e Considerações Adicionais}
\begin{itemize}
    \item Estudos observacionais e ensaios clínicos
    \item Relatórios de estudos em línguas não inglesas
    \item Literatura cinzenta (“grey literature”)
    \item Ano de publicação
    \item Refinamento de critérios ao longo da revisão
\end{itemize}

\begin{table}
\centering
\begin{tabular}{l l}
\toprule
Tipo de Estudo & Consideração \\
\midrule
ECR & Maior nível de evidência \\
Observacional & Pode fornecer dados suplementares \\
Relatórios não-ingleses & Considerar tradução e relevância \\
Literatura cinzenta & Inclui teses, relatórios técnicos \\
Ano de publicação & Definir período relevante \\
\bottomrule
\end{tabular}
\caption{Tipos de estudos e considerações}
\end{table}
\end{frame}

%-----------------------------------
\section{Exercício 1: Aplicando PICOTS}
\begin{frame}{Exercício 1: Aplicando PICOTS}
\begin{itemize}
    \item Questões básicas sobre cada elemento PICOTS
    \item Aplicação de critérios de elegibilidade
    \item Discussão: “o que faria com…?” (sem respostas certas)
\end{itemize}

\begin{figure}
\centering
\includegraphics[width=0.7\textwidth]{exercicio_picots.png} % Placeholder
\caption{Exercício prático de PICOTS}
\end{figure}
\end{frame}

%-----------------------------------
\section{Refinamento de Critérios de Elegibilidade}
\begin{frame}{Refinamento de Critérios de Elegibilidade}
\begin{itemize}
    \item Avaliação contínua dos critérios durante a revisão
    \item Ajuste de critérios amplos ou restritos conforme necessidade
    \item Minimização de viés
\end{itemize}
\end{frame}

%-----------------------------------
\section*{Referências}
\begin{frame}{Referências}
\begin{itemize}
    \item Higgins JPT, Thomas J, Chandler J, et al. \textit{Cochrane Handbook for Systematic Reviews of Interventions}, 2nd Edition, 2019.
    \item Borenstein M., Hedges L., Higgins J., Rothstein H. \textit{Introduction to Meta-Analysis}, 2009.
    \item Mucomole J., Silva A., Pereira B. (2025). \textit{Energy-related study in Mozambique}. \href{https://doi.org/10.3390/en18061460}{DOI: 10.3390/en18061460}
\end{itemize}
\end{frame}

\end{document}
